\documentclass[11pt]{article}
\usepackage[margin=1in]{geometry}
\usepackage{amsmath,amssymb,amsthm}
\usepackage[colorlinks=true,linkcolor=blue,citecolor=blue,urlcolor=blue]{hyperref}
\newtheorem{theorem}{Theorem}[section]
\newtheorem{definition}[theorem]{Definition}
\newtheorem{assumption}[theorem]{Assumption}
\theoremstyle{remark}
\newtheorem{remark}[theorem]{Remark}

\title{Modular Recovery from Split Inclusions\\
{\large A B-minimal bridge from QFT collar geometry to approximate state
reconstruction}\\
{\normalsize Version 2}}
\author{Lluis Eriksson\\ Independent Researcher\\ \texttt{lluiseriksson@gmail.com}}
\date{January 2026 (v1); July 2026 (v2)}

\begin{document}
\maketitle

\begin{abstract}
In algebraic quantum field theory (AQFT), local algebras are typically
Type III factors, so density matrices and von Neumann entropies are
unavailable for bounded regions. We formulate a B-minimal continuum analog
of the lattice ``collar $\Rightarrow$ Markovness $\Rightarrow$ recovery''
mechanism by combining: (i) the split property as the mathematical
replacement of a buffer (collar), (ii) Araki relative entropy to define a
split-regularized conditional mutual information $I^{\mathcal{N}}_\omega
(A:C|B)$ relative to fixed Type I interpolating data $\mathcal{N}$, and
(iii) modular/twirled Petz recovery as an explicit candidate recovery
channel. Assuming an FR-type recoverability inequality in the fixed-split
setting, we obtain quantitative recovery bounds in a fidelity-based error
metric (purified distance). We conclude with a conditional holographic
remark. \emph{v2 corrects the fidelity-convention factor in the assumed
FR-type inequality:} with the squared (Bures/Uhlmann-squared) convention
used throughout, the importable finite-dimensional motivation gives
$-\log F_B \le I$, not $-\log F_B \le \tfrac12 I$; v1's half-form is
strictly stronger than its motivation and is refuted as a
finite-dimensional anchor on 40/40 random tripartite states (the corrected
form holds on 40/40) -- the same factor-2 correction applied in 2512.0101
v2 and 2601.0020 v2. The recovery theorem's constant changes by $\sqrt2$;
the exponent is unaffected. v2 further adds series positioning, a
compatibility remark with the split-regularized CMI of 2601.0020, a
concrete demonstration that $I^{\mathcal{N}}$ can be negative for
unfavorable split data -- nonnegativity of $I^{\mathcal{N}}$ is part of
the good-split-data regime, not automatic -- and a finite-dimensional
verification suite for every checkable ingredient of the dictionary.
\end{abstract}

\section{Introduction}
In finite-dimensional quantum systems, small conditional mutual
information (CMI) implies approximate recoverability \cite{FR}. In
continuum QFT, local algebras are typically Type III \cite{Haag,
Yngvason}: no density matrices, no finite von Neumann entropies for
bounded regions. This note provides a minimal framework importing the
CMI$\Rightarrow$recovery mechanism into AQFT, with the dictionary:
collar/buffer $\rightsquigarrow$ split property; von Neumann entropy
$\rightsquigarrow$ Araki relative entropy; Petz/twirled Petz
$\rightsquigarrow$ modular (twirled) Petz channel. The CMI here is
split-regularized relative to fixed Type I interpolating data
$\mathcal{N}$, and the FR-type bound is an \emph{assumption}: the
contribution is a clean continuum dictionary separating geometric input
(control of $I^{\mathcal{N}}_\omega$) from information-theoretic input
(recoverability), not a universal Type III recovery theorem.

\paragraph{Series positioning (added in v2).} This note is the Type III
capstone of the clustering--recovery program: it is the target of the
conjectural sections of 2512.0060 (Conj.~5.1) \cite{S0060}, 2512.0101
(Conj.~1, CMI route, whose v2 fixes the same factor corrected here)
\cite{S0101}, 2601.0007 (Conj.~8.1 and its Type III roadmap via
noncommutative $L^p$) \cite{S0007}, and 2601.0031 (Conj.~6.4)
\cite{S0031}. The companion 2601.0020 \cite{S0020} defines a
split-regularized CMI by a complementary route; see
Remark~\ref{rem:compat}.

\section{Preliminaries}
We assume a local net $\mathcal{O} \mapsto \mathcal{A}(\mathcal{O})
\subset \mathcal{B}(\mathcal{H})$ with isotony and locality \cite{Haag,
BR}; for bounded regions the algebras are typically Type III$_1$
\cite{Yngvason}, so ``discarding $C$'' is restriction $\omega_{AB} :=
\omega_{ABC}|_{\mathcal{A}_{AB}}$. An inclusion $\mathcal{M}_1 \subset
\mathcal{M}_2$ is \emph{split} if there is a Type I factor $\mathcal{N}$
with $\mathcal{M}_1 \subset \mathcal{N} \subset \mathcal{M}_2$
\cite{DL, BW}; the interpolating factor is not unique, and we fix one
choice wherever split is invoked (superscript $\mathcal{N}$). Araki
relative entropy $S(\omega\|\varphi)$ is defined via relative modular
operators \cite{Araki1, Araki2}; $F_B(\omega, \varphi) \in [0,1]$ denotes
Bures fidelity in the \emph{squared} convention (in Type I it reduces to
Uhlmann fidelity squared, $\|\sqrt\rho\sqrt\sigma\|_1^2$), and the
purified distance is $P(\omega,\varphi) := \sqrt{1 - F_B(\omega,
\varphi)}$.

\section{Split-regularized conditional mutual information}
Fix a tripartition $A$--$B$--$C$ with collar $B$ and a normal faithful
$\omega \in S(\mathcal{A}_{ABC})$.
\begin{assumption}[Fixed split data $\mathcal{N}$]
\label{ass:split}
The collar geometry supports split inclusions inducing fixed
identifications $\mathcal{A}_{AB} \cong \mathcal{A}_A \bar\otimes
\mathcal{A}_B$ and $\mathcal{A}_{ABC} \cong \mathcal{A}_A \bar\otimes
\mathcal{A}_{BC}$; the totality of fixed choices is denoted $\mathcal{N}$.
\end{assumption}
\begin{definition}[Split-regularized CMI]
$I^{\mathcal{N}}_\omega(A:B) := S(\omega_{AB}\|\omega_A \otimes
\omega_B)$, $I^{\mathcal{N}}_\omega(A:BC) := S(\omega_{ABC}\|\omega_A
\otimes \omega_{BC})$, and
\[
I^{\mathcal{N}}_\omega(A:C|B) := I^{\mathcal{N}}_\omega(A:BC) -
I^{\mathcal{N}}_\omega(A:B).
\]
\end{definition}
\begin{remark}[Dependence on $\mathcal{N}$ and (non-)negativity;
sharpened in v2]
\label{rem:Ndep}
We do not claim canonicity in $\mathcal{N}$, nor $I^{\mathcal{N}}_\omega
(A:C|B) \ge 0$ for arbitrary fixed data. v2 makes this concrete: in the
finite-dimensional dictionary check (Section~\ref{sec:verif}), taking a
state with $A$ nearly uncorrelated from $BC$ and rotating only the
$A|B$-layer identification by a generic unitary yields
$I^{\mathcal{N}} < 0$ on 20/20 sampled identifications (range $[-0.46,
-0.18]$ against $I(A:BC) = 0.05$). Nonnegativity is thus a genuine
restriction on good split data, not a formality.
\end{remark}
\begin{remark}[Compatibility with 2601.0020 (added in v2)]
\label{rem:compat}
2601.0020 \cite{S0020} defines a split-regularized CMI by embedding
$\iota_w: \mathcal{A}^{(w)}_{ABC} \hookrightarrow \mathcal{A}(ABC)$ and
computing the Type I CMI of the induced state $\omega \circ \iota_w$; here
we fix split identifications and take a difference of Araki relative
entropies. In Type I with the natural factorization both reduce to the
standard CMI (chain rule; verified to machine precision in the suite), so
they are complementary regularizations of the same target; away from the
natural choice they differ, and each carries its own scheme dependence.
\end{remark}

\section{Recovery channels and the modular/twirled Petz candidate}
Channels are normal CP unital maps in Heisenberg picture, acting on states
via the predual. A recovery channel on the collar is
$\mathcal{R}_{B\to BC,*}: S(\mathcal{A}_B) \to S(\mathcal{A}_{BC})$, and
the recovered state is $\tilde\omega_{ABC} := (\mathrm{id}_A \otimes
\mathcal{R}_{B\to BC,*})(\omega_{AB})$ under the fixed identifications.
The twirled modular Petz channel is, schematically,
$\mathcal{R}^{\rm tw}_{B\to BC} = \int dt\, \beta_0(t)\, \sigma^{\varphi_B}_t
\circ \mathcal{R}^{\rm mod}_{B\to BC} \circ \sigma^{\varphi_{BC}}_{-t}$
with $\beta_0(t) = \tfrac{\pi}{4}\,{\rm sech}^2(\pi t/2)$ (the universal
weight of \cite{JRSWW}; it integrates to 1 -- verified in the suite); the
rigorous construction uses standard form and Connes cocycles \cite{JRSWW,
SutterTH}.

\begin{assumption}[FR-type recoverability in the fixed-split setting;
factor corrected in v2]
\label{ass:fr}
In the setting of Assumption~\ref{ass:split}, assume
$I^{\mathcal{N}}_\omega(A:C|B)$ is finite and nonnegative and that there
exists a recovery channel $\mathcal{R}_{B\to BC,*}$ (e.g.\ the twirled
modular Petz channel) such that
\[
-\log F_B\bigl(\omega_{ABC},\, \tilde\omega_{ABC}\bigr) \;\le\;
I^{\mathcal{N}}_\omega(A:C|B).
\]
\end{assumption}
\begin{remark}[Correction note]
\label{rem:frfix}
v1 assumed $-\log F_B \le \tfrac12 I^{\mathcal{N}}$. With $F_B$ in the
squared convention (Definition of Section~2), the importable
finite-dimensional motivation -- universal recovery \cite{FR, JRSWW} --
gives $I \ge -\log F$; the $\tfrac12$ belongs to the root-fidelity form
$I \ge -2\log f$. The v1 half-form is \emph{strictly stronger} than its
motivation: on random 3-qubit states with Petz recovery the suite finds
$-\log F \le I$ on 40/40 draws but $-\log F \le \tfrac12 I$ on 0/40.
This is the same factor-2 correction applied in 2512.0101 v2 and
2601.0020 v2; stating the assumption in the importable form keeps the
conjectural target honest.
\end{remark}

\begin{theorem}[From split-CMI to recovery error (conditional; constant
updated in v2)]
\label{thm:main}
Assume Assumptions~\ref{ass:split} and \ref{ass:fr}. If for collar width
$w$ one has $I^{\mathcal{N}}_\omega(A:C|B) \le K e^{-\alpha w}$, then
there exists a recovery channel with
\[
P\bigl(\omega_{ABC}, \tilde\omega_{ABC}\bigr) \;\le\; \sqrt{K}\,
e^{-\alpha w/2}.
\]
\end{theorem}
\begin{proof}
$1 - F_B \le -\log F_B \le I^{\mathcal{N}} \le K e^{-\alpha w}$; take
square roots. (v1's constant $\sqrt{K/2}$ corresponded to the stronger
half-form assumption.)
\end{proof}

\section{What is proved vs.\ what is assumed}
\emph{Defined:} the split-dependent quantity $I^{\mathcal{N}}_\omega
(A:C|B)$ via Araki relative entropy under fixed split data.
\emph{Assumed:} the FR-type inequality (Assumption~\ref{ass:fr}), now in
its importable form. \emph{Concluded:} collar-width decay of
$I^{\mathcal{N}}$ implies decay of the purified-distance recovery error.
The split choice $\mathcal{N}$ is part of the regularization scheme.

\section{Physical hooks}
Split is expected under nuclearity \cite{BW}; Reeh--Schlieder provides
cyclic/separating vectors for modular theory. In massive QFTs with
exponential clustering it is natural to expect $I^{\mathcal{N}}_\omega
(A:C|B) \le K e^{-\alpha w}$ for fixed-size $A$ and appropriate
geometries -- this is precisely the continuum target of the lattice
program: Gaussian/lattice evidence in 2512.0060 \cite{S0060} and
2601.0007 \cite{S0007}, non-Gaussian CMI route in 2512.0101 \cite{S0101},
Gibbs-state CMI inputs collected in 2601.0020 \cite{S0020}.

\section{Conditional holographic remark}
As in v1: under a bulk--boundary interface in which subregion duality is
an approximate recovery statement for code subalgebras, a boundary collar
$B$ with small $I^{\mathcal{N}}_\omega(A:C|B)$ provides a sufficient
condition for approximate recovery from the collar, hence conditionally
for quantitative entanglement wedge reconstruction.

\section{Verification suite (added in v2)}
\label{sec:verif}
\texttt{verification/2601-0034/verify\_2601\_0034.py} checks every
finite-dimensional ingredient of the dictionary: (1) with natural Type I
data, $I^{\mathcal{N}} = $ standard CMI (chain rule, machine precision);
(2) the $\mathcal{N}$-dependence demo of Remark~\ref{rem:Ndep}
($I^{\mathcal{N}} < 0$ on 20/20 rotated identifications for an
$A$-uncorrelated state); (3) $\int \beta_0 = 1$; (4) the
finite-dimensional anchor of Assumption~\ref{ass:fr}: corrected form
40/40, v1 half-form 0/40 (Remark~\ref{rem:frfix}); (5) the arithmetic of
Theorem~\ref{thm:main} with the corrected constant.

\section{Discussion}
This B-minimal note isolates a continuum version of the
collar$\Rightarrow$recovery mechanism in AQFT, with explicit split
dependence and a clean separation of geometric and information-theoretic
inputs -- now with the assumed input stated in the form its
finite-dimensional motivation can actually support. Nonnegativity of $I^{\mathcal{N}}$ is part of the
good-split-data regime, not automatic (Remark~\ref{rem:Ndep}). Open
problems: proving Assumption~\ref{ass:fr} as a theorem in the fixed-split
setting with explicit constants (the Haagerup $L^p$ route sketched in 2601.0007's
Type III roadmap), and understanding stability with respect to the split
choice $\mathcal{N}$.

\section*{Note added in v2}
One correction and three additions, all verified where checkable. (1)
\emph{FR factor:} Assumption~\ref{ass:fr} restated in the importable
squared-fidelity form $-\log F_B \le I^{\mathcal{N}}$ (v1's half-form is
refuted as a finite-dimensional anchor on 40/40 draws);
Theorem~\ref{thm:main}'s constant becomes $\sqrt{K}$. Third occurrence of
this factor in the series (2512.0101 v2, 2601.0020 v2). (2)
\emph{Compatibility:} Remark~\ref{rem:compat} relates the present
$I^{\mathcal{N}}$ to the split-regularized CMI of 2601.0020 (both reduce
to standard CMI for natural Type I data; verified). (3)
\emph{$\mathcal{N}$-dependence made concrete:} Remark~\ref{rem:Ndep} now
carries an explicit negative-$I^{\mathcal{N}}$ demonstration. (4) Series
positioning and verification suite added; v1 cited no companion papers.

\begin{thebibliography}{16}
\bibitem{FR} O. Fawzi and R. Renner, Commun. Math. Phys. \textbf{340}
(2015) 575--611.
\bibitem{Haag} R. Haag, \emph{Local Quantum Physics}, 2nd ed., Springer
(1996).
\bibitem{Yngvason} J. Yngvason, Rep. Math. Phys. \textbf{55} (2005)
135--147.
\bibitem{BR} O. Bratteli and D. W. Robinson, \emph{Operator Algebras and
Quantum Statistical Mechanics 1}, 2nd ed., Springer (1987).
\bibitem{Araki1} H. Araki, Publ. RIMS Kyoto Univ. \textbf{11} (1976)
809--833.
\bibitem{Araki2} H. Araki, Publ. RIMS Kyoto Univ. \textbf{13} (1977)
173--192.
\bibitem{DL} S. Doplicher and R. Longo, Invent. Math. \textbf{75} (1984)
493--536.
\bibitem{BW} D. Buchholz and E. H. Wichmann, Commun. Math. Phys.
\textbf{106} (1986) 321--344.
\bibitem{SutterTH} D. Sutter, M. Tomamichel, and A. W. Harrow, IEEE Trans.
Inf. Theory \textbf{62} (2016) 2907--2913.
\bibitem{JRSWW} M. Junge, R. Renner, D. Sutter, M. M. Wilde, and A.
Winter, Ann. Henri Poincar\'e \textbf{19} (2018) 2955--2978.
\bibitem{S0060} L. Eriksson, ai.viXra:2512.0060 (2025; v2 2026).
\bibitem{S0101} L. Eriksson, ai.viXra:2512.0101 (2025; v2 2026).
\bibitem{S0007} L. Eriksson, ai.viXra:2601.0007 (2026; v2 2026).
\bibitem{S0020} L. Eriksson, ai.viXra:2601.0020 (2026; v2 2026).
\bibitem{S0031} L. Eriksson, ai.viXra:2601.0031 (2026; v2 2026).
\end{thebibliography}
\end{document}
